% Reduced extract from Vincent Goulet, Rédaction avec LaTeX.
% License: CC BY-SA 4.0. See SOURCES.md for attribution and source revision.
\documentclass{book}
\usepackage{booktabs,tabularx}

\begin{document}
\chapter{Tableaux et figures}
\label{chap:tableaux}

Les tableaux et graphiques ne sont pas les éléments de texte les plus simples
et rapides à créer avec \LaTeX. Pour ce type de contenu comme pour tout autre,
\LaTeX fait ce qu'on lui demande, sans tenter de deviner notre pensée.

\section{De la conception de beaux tableaux}
\label{sec:tableaux:booktabs}

Les tableaux servent à disposer de l'information sous forme de grille. Le
paquetage \texttt{booktabs} recommande des filets horizontaux bien choisis et
évite les filets verticaux. Voici un extrait de tableau adapté du manuel:

\begin{table}
\caption{Un exemple de tableau aéré}\label{tab:beau}
\begin{tabularx}{\textwidth}{>{\bfseries}lXr}
\toprule
Etape & Description & Valeur \\
\midrule
0 & Le premier calcul reste lisible. & 52 \\
1 & La largeur de cette colonne est ajustable. & 24 \\
\bottomrule
\end{tabularx}
\end{table}

La seconde version est plus aérée et plus facile à consulter. Comme expliqué
dans la section \autoref{sec:tableaux:booktabs}, les tableaux clairs ne
nécessitent pas de filets verticaux.
\end{document}
